\chapter{Einleitung}\label{sec:einleitung}

Supraleitung zählt zu den faszinierendsten makroskopischen Quantenphänomenen der Festkörperphysik: Unterhalb einer kritischen Temperatur verschwindet der elektrische Widerstand eines Supraleiters vollständig und Magnetfelder werden vollständig aus seinem Inneren verdrängt (Meißner-Ochsenfeld-Effekt). 
Diese Eigenschaften ermöglichen Anwendungen wie die verlustfreie Energieübertragung, hochauflösende Magnetresonanztomographie oder die Realisierung von Qubits in Quantencomputern. 
Die Untersuchung der supraleitenden Eigenschaften von Materialien und der ihnen zugrunde liegenden Mechanismen ist daher sowohl aus technologischer als auch aus grundlagenphysikalischer Sicht von großem Interesse.

Das Bienenwabengitter stellt am Beispiel von Graphen eine Klasse zweidimensionaler
Materialien mit einer Reihe außergewöhnlicher elektronischer Eigenschaften dar. 
Aufgrund seiner linearen Bandstruktur in der Nähe der Dirac-Punkte treten in Graphen unter anderem masselose Dirac-Fermionen~\cite{semenoff_condensed-matter_1984,novoselov_two-dimensional_2005} sowie eine am Dirac-Punkt linear verschwindende Zustandsdichte auf \cite{castro_neto_electronic_2009}. 
Reines, undotiertes Graphen zeigt trotz dieser besonderen Eigenschaften jedoch keine Supraleitung. 
Bislang konnte diese experimentell nur in verdrehtem, dopellagigem Graphen~\cite{cao_unconventional_2018} sowie in dotiertem, einlagigem Graphen~\cite{ludbrook_evidence_2015} direkt nachgewiesen werden. 
Diese Beobachtungen werfen die Frage auf, unter welchen Bedingungen sich in Materialien mit Bienenwabenstruktur überhaupt eine supraleitende Phase ausbilden kann, und ob dabei auch der konventionelle, phononenvermittelte Mechanismus der BCS-Theorie eine Rolle spielen könnte.

Ziel dieser Arbeit ist es daher, die durch die Gitterstruktur bedingten Eigenschaften konventioneller Supraleitung auf dem Bienenwabengitter zu untersuchen, um für diese Materialklasse und insbesondere für Graphen genauere Vorhersagen treffen zu können, unter welchen Bedingungen Supraleitung zu erwarten ist.

Dazu werden in Kapitel~\ref{sec:grundlagen} zunächst die theoretischen Grundlagen zur Geometrie und elektronischen Struktur des Bienenwabengitters sowie zur BCS-Theorie eingeführt.  
In Kapitel~\ref{sec:modell} wird darauf aufbauend ein Nächste-Nachbar-Tight-Binding-Modell mit BCS-Wechselwirkung formuliert und mittels Mean-Field-Näherung eine Selbstkonsistenzgleichung für den Ordnungsparameter $\Delta$ hergeleitet, sowohl für den Spezialfall halber Füllung als auch für beliebige Füllung.
Kapitel~\ref{sec:ergebnisse} präsentiert die numerische Lösung dieser Gleichungen:
Zunächst werden die allgemeinen Abhängigkeiten von $\Delta$ und der kritischen Temperatur $T_C$ von den Modellparametern untersucht, unter anderem ein quantenkritischer Punkt bei halber Füllung. 
Anschließend wird das Modell auf realistische Parameter für Graphen angewendet und mit Vorhersagen aus der Literatur verglichen. 
Kapitel~\ref{sec:fazit} fasst die Ergebnisse zusammen und gibt einen
Ausblick auf mögliche Erweiterungen des Modells.